2. Plantee las siguientes integrales dobles en ambos órdenes de integración. Después resuelva la integral usando el orden más sencillo (SÓLO resuelva uno).\\
a) \( \iint_{D} y d A \) donde \( D \subset \mathbb{R}^{2} \) está acotada por \( y=x-2 \) y \( x=y^{2} \). \\
La región \( D \) está delimitada por \( y = x - 2 \) y \( x = y^2 \). Esto nos da los siguientes despejes bastante simples:
\begin{itemize}
    \item Región tipo I: \( 2 \leq x \leq 3 \), \( \sqrt{x} \leq y \leq x - 2 \).  
    \item Región tipo II: \( 0 \leq y \leq 1 \), \( y^2 \leq x \leq y + 2 \).  
\end{itemize}

Ambas parecen casi igual de difíciles, pero en general las raíces a veces son complicadas de interpretar así que resolvemos en Tipo II:  

\[
I = \int_0^1 \int_{y^2}^{y+2} y \, dx \, dy = \int_0^1 y (y + 2 - y^2) dy = \int_0^1 (y^2 + 2y - y^3) dy = \left( \frac{1}{3} + 1 - \frac{1}{4} \right) = \frac{13}{12}.
\]

b) \( \iint_{D} y^{2} e^{x y} d A \) donde \( D \subset \mathbb{R}^{2} \) está acotada por \( y=x, y=4 \) y \( x=0 \). \\
La región \( D \) está acotada por \( y = x \), \( y = 4 \) y \( x = 0 \). Este es más sencillo que el anterior pues tenemos y=x así que solo cambia una pequeña parte.
\begin{itemize}
    \item Región tipo I: \( 0 \leq x \leq 4 \), \( x \leq y \leq 4 \).  
    \item Región tipo II: \( 0 \leq y \leq 4 \), \( 0 \leq x \leq y \).  
\end{itemize}

Resolvemos en Tipo II:  

\[
I = \int_0^4 \int_0^y y^2 e^{xy} \, dx \, dy = \int_0^4 y^2 \left[ \frac{e^{xy}}{y} \Big|_0^y \right] dy = \int_0^4 y^2 \frac{e^{y^2} - 1}{y} dy = \int_0^4 (y e^{y^2} - y) dy
\]
\[
= \frac{1}{2} \int_0^{16} e^u du - \int_0^4 y dy. = \frac{e^{16} - 1}{2} - 8.
\]  

c) \( \iint_{D} \operatorname{sen}(x) d A \) donde \( D=\left\{(x, y) \in \mathbb{R}^{2} ; \frac{-\pi}{2} \leq y \leq \frac{\pi}{2},-y-\frac{\pi}{2} \leq x \leq y\right\} \). \\

La región \( D \) está definida por \( -\frac{\pi}{2} \leq y \leq \frac{\pi}{2} \) y \( -y - \frac{\pi}{2} \leq x \leq y \). De aquí nótese que tenemos.
\begin{itemize}
    \item Región tipo I: \( -\frac{\pi}{2} \leq x \leq \frac{\pi}{2} \), \( -x - \frac{\pi}{2} \leq y \leq x \).  
    \item Región tipo II: \( -\frac{\pi}{2} \leq y \leq \frac{\pi}{2} \), \( -y - \frac{\pi}{2} \leq x \leq y \).  
\end{itemize}

Resolvemos en Tipo II:  

\[
I = \int_{-\frac{\pi}{2}}^{\frac{\pi}{2}} \int_{-y - \frac{\pi}{2}}^y \sin(x) \, dx \, dy.
\]

Resolvemos primero la integral interna que es bastante simple,

\[
\int \sin(x) dx = -\cos(x).
\]

Evaluamos en los límites

\[
-\cos(y) + \cos(-y - \frac{\pi}{2}).
\]

Como \( \cos(-y - \frac{\pi}{2}) = \sin(y) \), tenemos:

\[
I = \int_{-\frac{\pi}{2}}^{\frac{\pi}{2}} (-\cos(y) + \sin(y)) dy = -\int_{-\frac{\pi}{2}}^{\frac{\pi}{2}} \cos(y) dy + \int_{-\frac{\pi}{2}}^{\frac{\pi}{2}} \sin(y) dy -\sin(y) \Big|_{-\frac{\pi}{2}}^{\frac{\pi}{2}} + (-\cos(y)) \Big|_{-\frac{\pi}{2}}^{\frac{\pi}{2}} 
\]
\[
= - (1 - (-1)) + (-0 + 0) = - (1 + 1) + 0 = -2 = -2
\]

d) \( \iint_{D} \frac{y}{x^{5}+1} d A \) donde \( D=\left\{(x, y) \in \mathbb{R}^{2} ; 0 \leq x \leq 1,0 \leq y \leq x^{2}\right\} \). \\
La región \( D \) está dada por \( 0 \leq x \leq 1 \), \( 0 \leq y \leq x^2 \). Se sigue que.
\begin{itemize}
    \item Región tipo I: \( 0 \leq x \leq 1 \), \( 0 \leq y \leq x^2 \).  
    \item Región tipo II: \( 0 \leq y \leq 1 \), \( 0 \leq x \leq \sqrt{y} \).  

Resolvemos en Tipo I:  

\[
I = \int_0^1 \int_0^{x^2} \frac{y}{x^5+1} dy \, dx.
\]

Hacemos la integral interna como de costumbre:

\[
\int_0^{x^2} \frac{y}{x^5+1} dy = \frac{1}{x^5+1} \int_0^{x^2} y dy = \frac{1}{x^5+1} \cdot \frac{y^2}{2} \Big|_0^{x^2} = \frac{x^4}{2(x^5+1)}.
\]

Ahora la integral externa:

\[
I = \int_0^1 \frac{x^4}{2(x^5+1)} dx.
\]

Es demasiado sugerente el logaritmo natural como primitiva de la función integrada, pero lo haremos más visible con el cambio \( u = x^5 + 1 \), con \( du = 5x^4 dx \), veamos que entonces:

\[
I = \frac{1}{10} \int_1^2 \frac{du}{u} = \frac{1}{10} \ln |u| \Big|_1^2 = \frac{1}{10} (\ln 2 - \ln 1) = \frac{\ln 2}{10}.
\]